% ============================================================
% Template LaTeX ABNT NBR 14724/2024 — IFB
% Compatível com pdflatex + abntex2cite
% ============================================================

\documentclass[
  12pt,
  a4paper,
  oneside,
  brazil
]{article}

% ── Codificação e Idioma ──────────────────────────────────────────────────────
\usepackage[T1]{fontenc}
\usepackage[utf8]{inputenc}
\usepackage[brazil]{babel}

% ── Fonte Times New Roman ─────────────────────────────────────────────────────
\usepackage{mathptmx}

% ── Geometria da página — margens ABNT anverso ────────────────────────────────
\usepackage[
  a4paper,
  left=3cm,
  top=3cm,
  right=2cm,
  bottom=2cm,
  headheight=14pt,
  headsep=0.5cm
]{geometry}

% ── Espaçamento 1,5 entre linhas (ABNT) ──────────────────────────────────────
\usepackage{setspace}
\onehalfspacing

% ── Recuo do parágrafo: 1,25cm (ABNT) ────────────────────────────────────────
\usepackage{indentfirst}
\setlength{\parindent}{1.25cm}
\setlength{\parskip}{0pt}

% ── Formatação de seções — numeração progressiva NBR 6024 ────────────────────
\usepackage{titlesec}

\titleformat{\section}
  {\normalfont\normalsize\bfseries\MakeUppercase}
  {\thesection\quad}
  {0em}{}

\titleformat{\subsection}
  {\normalfont\normalsize\bfseries}
  {\thesubsection\quad}
  {0em}{}

\titleformat{\subsubsection}
  {\normalfont\normalsize\bfseries\itshape}
  {\thesubsubsection\quad}
  {0em}{}

\titleformat{\paragraph}[runin]
  {\normalfont\normalsize\itshape}
  {\theparagraph\quad}
  {0em}{}[\hspace{0.5em}]

\titleformat{\subparagraph}[runin]
  {\normalfont\normalsize}
  {\thesubparagraph\quad}
  {0em}{}[\hspace{0.5em}]

\titlespacing*{\section}      {0pt}{18pt}{12pt}
\titlespacing*{\subsection}   {0pt}{18pt}{6pt}
\titlespacing*{\subsubsection}{0pt}{12pt}{6pt}
\titlespacing*{\paragraph}    {0pt}{12pt}{0pt}

% Seções primárias iniciam em nova página
\newcommand{\secbreak}{\clearpage}
\let\oldsection\section
\renewcommand{\section}{\secbreak\oldsection}

% ── Numeração de páginas — canto superior direito, 10pt ──────────────────────
\usepackage{fancyhdr}
\pagestyle{fancy}
\fancyhf{}
\fancyhead[R]{\fontsize{10}{12}\selectfont\thepage}
\renewcommand{\headrulewidth}{0pt}
\renewcommand{\footrulewidth}{0pt}

\fancypagestyle{plain}{
  \fancyhf{}
  \fancyhead[R]{\fontsize{10}{12}\selectfont\thepage}
  \renewcommand{\headrulewidth}{0pt}
}

% ── Tabelas ───────────────────────────────────────────────────────────────────
\usepackage{longtable}
\usepackage{booktabs}
\usepackage{array}
\usepackage{tabularx}
\renewcommand{\arraystretch}{1.3}

% ── Legendas ─────────────────────────────────────────────────────────────────
\usepackage{caption}
\DeclareCaptionFormat{abnt}{\fontsize{10}{12}\selectfont#1#2#3}
\captionsetup{
  format=abnt,
  justification=centering,
  singlelinecheck=false,
  labelsep=emdash,
  position=above,
  skip=3pt
}

\newcommand{\fonte}[1]{%
  \begin{flushleft}
    \fontsize{10}{12}\selectfont Fonte: #1
  \end{flushleft}
}

% ── Listas ───────────────────────────────────────────────────────────────────
\usepackage{enumitem}
\setlist{nosep, leftmargin=*}
\def\tightlist{}

% ── Sumário ──────────────────────────────────────────────────────────────────
\usepackage{tocloft}
\renewcommand{\contentsname}{SUMÁRIO}

% ── Hyperlinks ───────────────────────────────────────────────────────────────
\PassOptionsToPackage{hyphens}{url}
\usepackage[hidelinks, colorlinks=false, breaklinks=true]{hyperref}

% ── Citações longas ──────────────────────────────────────────────────────────
\usepackage{quoting}
\quotingsetup{font=small, indentfirst=false, vskip=\baselineskip}

% ── Figuras ──────────────────────────────────────────────────────────────────
\usepackage{graphicx}
\usepackage{float}

% ── Referências ABNT ─────────────────────────────────────────────────────────
\usepackage[
  alf,
  abnt-emphasize=bf,
  abnt-etal-cite=2,
  abnt-etal-list=0
]{abntex2cite}

\date{}

% ============================================================
\begin{document}

% ══════════════════════════════════════════════════════════════════════════════
% CAPA — NBR 14724/2024
% ══════════════════════════════════════════════════════════════════════════════
\begin{titlepage}
  \thispagestyle{empty}
  \begin{center}

    \includegraphics[width=0.4\textwidth]{ifb_vs_dp.png}\\[1.8em]

    {\normalsize\bfseries
      Instituto Federal de Educação, Ciência e Tecnologia de Brasília --- IFB\\[0.3em]
      \textit{Campus} Samambaia\\[0.3em]
      Curso Superior de Tecnologia em Design de Produto
    }

    \vfill

    {\large\bfseries
      DEFINIÇÃO DO PROBLEMA --- MESA DE APOIO PILOTO
    }

    \vfill

    {\normalsize Brasília/DF\\2026}

  \end{center}
\end{titlepage}
\setcounter{page}{2}

% ══════════════════════════════════════════════════════════════════════════════
% FOLHA DE ROSTO — NBR 14724/2024
% ══════════════════════════════════════════════════════════════════════════════
\begin{titlepage}
  \thispagestyle{empty}
  \begin{center}

    {\normalsize\bfseries
      VICTOR COSTA\\
      ANNA JI\\
      LUÍS COSTA
    }

    \vfill

    {\large\bfseries DEFINIÇÃO DO PROBLEMA --- MESA DE APOIO PILOTO}

    \vfill

    \begin{flushright}
      \begin{minipage}{0.5\textwidth}
        \normalsize
        Relatório de etapa de projeto apresentado ao Curso Superior de
        Tecnologia em Design de Produto do Instituto Federal de Educação,
        Ciência e Tecnologia de Brasília --- IFB, \textit{Campus} Samambaia,
        como requisito parcial para aprovação na disciplina de Metodologia
        e Prática de Projeto.

        \vspace{1em}
        Professora: Profª. Drª. Fernanda Torres
      \end{minipage}
    \end{flushright}

    \vfill

    {\normalsize Brasília/DF --- 2026}

  \end{center}
\end{titlepage}

% ══════════════════════════════════════════════════════════════════════════════
% RESUMO — NBR 6028/2003
% ══════════════════════════════════════════════════════════════════════════════
\clearpage
\thispagestyle{empty}
\begin{center}
  {\normalsize\bfseries RESUMO}
\end{center}

\vspace{1em}
\noindent
O presente documento registra a etapa de definição do problema do projeto
Mesa de Apoio Piloto, desenvolvido no âmbito da disciplina de Metodologia
e Prática de Projeto. São apresentados o escopo do produto, suas funções,
o público-alvo e as referências estéticas que orientarão o desenvolvimento
projetual. O produto consiste em uma mesa de apoio/cabeceira de design autoral
com identidade estética fundamentada na arquitetura modernista brasileira,
com linguagem formal inspirada em Brasília, voltada a adultos das classes A e B
que valorizam mobiliário com referência cultural. As diretrizes formais,
cromáticas e materiais apresentadas fundamentam-se no vocabulário visual
da cidade e do modernismo nacional, com destaque para o freijó como material
de interesse e a construção modular por encaixe como abordagem construtiva
a ser investigada.

\vspace{1em}
\noindent\textbf{Palavras-chave:} design de produto; mobiliário autoral;
modernismo brasileiro; definição de problema; Brasília.

% ══════════════════════════════════════════════════════════════════════════════
% SUMÁRIO — NBR 14724/2024
% ══════════════════════════════════════════════════════════════════════════════
\clearpage
\thispagestyle{empty}
\tableofcontents

% ══════════════════════════════════════════════════════════════════════════════
% CONTEÚDO
% ══════════════════════════════════════════════════════════════════════════════

\section{Definição do Problema}

\subsection{O que?}

O problema de projeto consiste em desenvolver uma mesa de cabeceira / mesa
de apoio de sala de estar de design autoral com identidade estética
fundamentada na arquitetura modernista brasileira. A proposta exige uma
estrutura modular em madeira com painel frontal trabalhado e apoios laterais
escultóricos, de modo que o móvel cumpra sua função de suporte e
armazenamento sem abrir mão de uma presença visual própria no ambiente.
O objetivo é conciliar funcionalidade técnica, referência cultural explícita
e viabilidade produtiva, utilizando materiais disponíveis no mercado regional.

\subsection{Para que?}

A mesa de apoio deve atender aos seguintes requisitos:

\begin{itemize}
  \item \textbf{Dimensões versáteis}: Deve atender dentro das margens
    ergonômicas tanto o uso como mesa de cabeceira (quarto), quanto mesa
    de apoio (sala de estar);
  \item \textbf{Armazenamento e organização de objetos decorativos}: prateleiras
    internas/externas e/ou módulos internos para acomodar pequenos objetos
    variados de acordo com o contexto de uso;
\end{itemize}

\subsection{Para quem?}

O público-alvo são adultos entre 30 e 50 anos, das classes A e B, moradores
de apartamentos ou casas com arquitetura contemporânea ou minimalista.
São pessoas que já realizaram as principais aquisições do mobiliário doméstico
e passam a buscar substituição por peças de maior identidade: consumidores
que valorizam design autoral brasileiro, investem em objetos com referência
cultural e buscam diferenciação em relação ao mobiliário industrializado em
série disponível nas grandes redes.

Esse público reconhece referências da arquitetura e do modernismo nacional
e, quando encontra uma peça com história e autoria, tende a pagar mais por ela.

\subsection{Referências Estéticas}

\subsubsection{Formas}

Três matrizes formais orientam o projeto.

A geometrização marcante e a silhueta baixa seguem a tradição do mobiliário
\textit{mid-century} e se alinham com a arquitetura de Brasília. Quando
pensada em composição com um sofá ou uma cama, a proporção horizontalizada
e a altura reduzida liberam a parede como plano de fundo, traçando um paralelo
com o céu amplamente visível por conta dos prédios residenciais da cidade,
limitados a seis andares.

Os perfis em arco côncavo têm referência direta nas estruturas de Oscar
Niemeyer: os pilotis curvos, as arcadas do Palácio do Planalto, a estrutura
do Congresso Nacional. A curvatura distribui a carga para as extremidades e
cria uma tensão visual entre a massa horizontal da composição (objeto de apoio
+ objeto principal associado) e os apoios, que devem conversar com o conceito
de pilotis.

Padrões geométricos e modulação completam a linguagem formal, remetendo à
racionalidade urbana de Brasília: o traçado do plano piloto, o ritmo das
fachadas da Esplanada dos Ministérios, as tramas em azulejo do modernismo
brasileiro.

\subsubsection{Cores e Materiais}

A paleta deve nascer do material, não de uma decisão de cor: madeira natural
em tons âmbar e mel no plano dominante, com o contraste surgindo naturalmente
do recuo dos recortes e interiores, que ficam em sombra, sem pigmento nem
acabamento opaco. O uso de concreto, mármore e metal pode surgir como opção,
mas o foco é a madeira.

O freijó (\textit{Cordia goeldiana}) é o material de maior interesse nesse
sentido, uma madeira brasileira de tom dourado-âmbar amplamente disponível no
mercado nacional e com boa trabalhabilidade para usinagem. Sua associação ao
mobiliário autoral nacional reforça a identidade cultural da peça sem recorrer
a espécies de extração predatória; a escolha depende ainda de análise mais
profunda sobre custos e viabilidade dentro do contexto acadêmico.

A construção por encaixe é uma abordagem que seria de interesse ao projeto,
pois favorece a produção em marcenaria de pequeno porte, reduz desperdício de
material e permite transporte desmontado. Solução logística compatível com o
perfil de comercialização de móveis autorais e, não por coincidência, a mesma
lógica construtiva que estruturas modulares em madeira já provaram funcionar
muito bem.

\end{document}
